\documentclass{article}
\usepackage{mathtools}
\usepackage{graphicx}
\usepackage{xcolor}
\usepackage{url}
\usepackage[margin=1in]{geometry}

\begin{document}
\title{Blackhat Video Review - BLERP}
\author{Dustin Randall}
\date{}
\maketitle

\section{Summary}
I scrolled through the available Black Hat presentations and when I saw the word BLERP, I was hooked.
I realized that it was a talk about Bluetooth Low Energy, which is being used in one of the senior projects this year.
The video can be found at \url{https://www.youtube.com/watch?v=_08bb55Q33o}.

A grad student and his advisor presented the vulnerabilities they discovered in the BLE re-pairing process.
This work was important because BLE runs on billions of devices and any vulnerabilities in the standard can leave those devices vulnerable.
The presentation started by explaining the basic pairing process in BLE and focused on the SL or Security Level, the PK or Pairing Key, and the E or Entropy level.
The standard for BLE is over 3800 pages, and there are only 4 references to "re-pairing" when grepped.

Overall, they discovered 6 separate vulnerabilities in the re-pairing process and developed 4 separate attacks.
These vulnerabilities are all attacking the fact that during re-pairing the E level is not verified, and neither is the new Security Level.
This enables a downgrade attack where the attacker can reduce the entropy level to something which can be cracked by brute force.
The four attacks allow impersonating either the central or peripheral device, as well as performing 2 different man-in-the-middle attacks.

They tested their attacks against 22 different devices, including 13 vendors, 16 central devices, and 9 peripheral devices.
Nearly all the devices were vulnerable to all the attacks with Windows and Linux PCs being partially resistant.
During the talk, they showed a demo where a phone was connected to a mouse, and they showed that a re-pairing attack allowed a python script to connect to the phone impersonating the mouse.
Another slide showed that they had created a stand-alone device which could perform the attack, and have open-sourced the code to run it from a Linux machine.

The researchers did the responsible disclosure of their findings and attained CVE-2025-62235.
They also proposed 2 primary mitigations to the standard which would prevent all 6 vulnerabilities.
Changing standards is difficult and time-consuming, so they also provided mitigations which vendors could apply without departing from the existing standard.
When they provided these findings to the vendors, the majority of them did not implement the suggested mitigations, leaving many devices still vulnerable.
As of 2026 the standard remains unchanged as well.

\section{Presentation}
I think that the presentation was put together well.
The slides were clear and visual without being overly textual.
I felt that the progression was logical and well-paced, and that I did not require a lot of background knowledge to follow their main points.
The one minor issue I noticed is that the grad student tended to hold the microphone away from his face, which made the audio hard to hear at times.

\section{Content}
I really enjoyed the content of the presentation.
It certainly highlighted the vulnerability in so many devices, and the reality of politics in industry.
I'm not sure what the potential for impact is with performing these attacks with BLE devices in the wild.
They suggested maybe smartwatches could access sensitive data like health information.
I think that central devices must already be skeptical of peripheral devices which may mitigate the level of risk posed by impersonating a peripheral device.

\end{document}